\begin{document}
\docTitle{Guide to the notes system}
This document describes the goals and architecture of the notes system used for daily note-taking and knowledge system.

\section*{Design brief}
The goal of this system is to provide a structured workflow for writing, reading, reviewing, and preserving notes.
The system is based on the following setup:
\begin{itemize}
  \item notes are written and edited only on a computer;
  \item notes can be read, reviewed, and commented also on a phone;
\end{itemize}
The system follows these principles:
\begin{itemize}
\item LaTeX is used as the main format for writing notes;
\item comments are stored as one-line JSONL mappings in \code{comments/<pdf-stem>_comments.jsonl};
Comments are either updated on saving \code{current_pdfs/} or when running \code{tools/extract_comments.py}.
\item the phone and computer are automatically synchronized, so no manual action is required beyond writing notes and adding comments.
\end{itemize}

\subsection*{Phone workflow}
\subsubsection*{Android App}
PDFs are searched and explored in a smartphone app that lists each note.
Each entry shows the following information about each note:
\begin{itemize}
  \item when the PDF was last modified;
  \item when the last comment was added;
  \item the total number of comments that need to be addressed.
\end{itemize}
When the user clicks on a note, it should open in Adobe Acrobat Reader, where the user can read documents comfortably and add comments using sticky notes or highlighted text.
In the background, notes are synchronized with the computer using Syncthing.

\subsection*{Computer Workflow}
On the computer, the user edits \code{.tex} files.
The computer workflow should support:
\begin{itemize}
  \item editing \code{.tex} source files;
  \item automatically rebuilding PDFs every N minutes;
  \item adding new \code{.tex} files and making the resulting PDFs available to the phone;
  \item watching \code{current_pdfs/} for Adobe/Preview saves and updating the corresponding \code{comments/*.jsonl};
  \item preserving previous comments by re-adding unresolved comments after rebuilding PDFs;
  \item automatically committing changes to Git.
\end{itemize}


\section*{Overview}
Notes are grouped by topic in folders.
Each \code{docs-*} project is a specific topic. 
The folder structure is explained in \hyperref[sec:doc-folder-structure]{Doc folder structure}.

Additionally, there are additional top folders with shared utilities or where the pds are located. 
These are \code{current_pdfs/}, \code{incoming_pdfs/}, \code{library/},  \code{tools/}
The \code{tools/} folder contains tools for building PDFs and related automation, including a default latex package \code{manual_style.sty} which should be imported by each \code{.tex} fille.

This defines additional elements such as \code{code} and \code{DocCode} blocks.


\section*{Comments - Technical implementation}
The canonical comment store is the root \code{comments/} folder.
For each current PDF \code{current_pdfs/<pdf-stem>.pdf}, comments are stored in:
\begin{DocCode}
  comments/<pdf-stem>_comments.jsonl
\end{DocCode}
The merge baseline used for retaining comments across rebuilds is stored in:
\begin{DocCode}
  comments/.base/<pdf-stem>_comments.jsonl
\end{DocCode}

\subsection*{Comment merge model}
Comment synchronization uses three states:
\begin{itemize}
  \item base state: \code{comments/.base/<pdf-stem>_comments.jsonl}, the last synchronized state;
  \item comment state: \code{comments/<pdf-stem>_comments.jsonl}, the current machine-readable state;
  \item PDF state: annotations extracted from the target PDF, temporarily stored as \code{comments/tmp.jsonl}.
\end{itemize}
The merge compares the comment state and PDF state against the base state.
\begin{center}
\begin{tabular}{lll}
Comment state & PDF state & Result \\
unchanged & unchanged & keep either state \\
changed & unchanged & use comment state \\
unchanged & changed & use PDF state \\
changed same way & changed same way & keep either state \\
changed differently & changed differently & use the newest \code{modified} timestamp \\
deleted & unchanged & delete \\
deleted & changed & keep the changed side \\
both deleted & both deleted & delete
\end{tabular}
\end{center}
File modification times are only used as a signal that extraction may be needed.
Correctness comes from record-level comparison against the base state.

PDF annotations are exported to one-line JSONL mappings under \code{`comments/`}.
Use \code{python3 tools/extract_comments.py} to extract comments from PDFs.
Specifically this script extracts comments from PDFs in \code{incoming_pdfs/} to \code{comments/}.
PDFs whose filenames contain \code{.sync-conflict} are treated as variants of the corresponding original PDF, and their extracted comments are combined into the same JSONL output file.

Use \code{python3 tools/add_comments.py} to add comments to the PDFS.
Before writing comments into a target PDF, this script should compare the base state, comment state, and PDF state, then write the newest merged state.
The schema for the JSONL comment records is described in \hyperref[app:comment-jsonl-schema]{Appendix: Comment JSONL schema}.


\phantomsection
\section*{Doc folder structure}\label{sec:doc-folder-structure}
Each doc should have its own git repository and be matched in the doc library.
\begin{DocCode}
  docs-agent/
  docs-ideas/
  docs-process/
  docs-template/
  docs-tsENV/
  library/
    README.md
  tools/
    render_all.sh
    render_loop.sh
    pull_update_submodules.sh
    push_submodule_changes.sh
    extract_comments.py
    add_comments.py
    manual_style.sty
    watch_current_pdfs_comments.py
  current_pdfs/
    agent.pdf
    \dots
    library.pdf
  incoming_pdfs/
  comments/
\end{DocCode}


\begin{DocCode}
docs-my-topic/
|-- .git/
|-- manual/
|   |-- docs_order_manifest.txt
|   |-- src/
|   |   |-- first_note.tex
|   |   `-- second_note.tex
|   |-- project_specific_manual_style.sty [OPTIONAL]
|-- pdf/
|   |-- my-topic_first_note.pdf
|   `-- my-topic_combined.pdf
|-- third_party/ [OPTIONAL]
|   |-- first.pdf
|   `-- second.pdf
`-- build/ [.gitignored]
\end{DocCode}
\code{docs_order_manifest.txt} defines the order in which files from \code{manual/src/} are included when building the combined PDF.
Each non-comment line should name one \code{.tex} file. 
Blank lines and comment lines are ignored.
If a source file is added or removed, update \code{docs_order_manifest.txt} intentionally so the combined PDF only includes the listed renderable files.
Thus, the source files in \code{manual/src/} is the actual truth; the manifest only controls combined-PDF ordering.
The file \code{*_combined.pdf} contains all the .tex plus eventual .tex that are also hyperlinked
Each \code{docs-*} project is a separate document repository with its own \code{manual/}, \code{pdf/}, and \code{build/} folders.
\code{project_specific_manual_style.sty and \code{third_party/} is optional



The same document repository should be listed in \code{library/README.md}.
\agent{Authored files whose basename begins with \code{README}, case-insensitively, are not allowed anywhere in a \code{docs-*} repository; \path{library/README.md} is the sole README entrypoint for the documentation workspace on GitHub.}
\begin{DocCode}
- [docs-my-topic](https://github.com/<user>/docs-my-topic) - Short description of the document set.
\end{DocCode}


\section*{Citations}
The standard format when doing a related work, literature review is the following (contained in \code{tools/manual_style.sty})
\begin{DocCode}
\newcommand{\relatedWorkMeta}[5]{%
  \leavevmode\par\nobreak\vspace{0.05em}%
  \noindent{\scriptsize\setlength{\parskip}{0pt}%
    Citation: #1.
    First online: #2.
    Venue: #3.
    Repository: #4
    AuthorsFromZH: #5
    \\}%
}
\end{DocCode}
For example:
\begin{DocCode}
\relatedWorkMeta
  {23* (queried June 2, 2026)}
  {Jun 2023}
  {NeurIPS 2023 Datasets and Benchmarks Track}
  {Official repository, \url{https://github.com/Jwoo5/ecg-qa}}
  {-}
\end{DocCode}
\code{AuthorsFromZH} indicates whether one of the authors is from Zurich and, if so, which one.
How this information is obtained is explained in \hyperref[app:related-work-information-method]{Appendix: Information collection method}.
Ciitations, documents that expect a different citation style needs to have ale where they explain it 

\section*{Shared policy guideline}
If a `.tex` document from another project is useful as supporting material, it may be referenced with a hyperlink.
When creating the combined version, the PDF generated from that referenced `.tex` document should also be appended at the end of the combined PDF.


\section*{Compiling into PDF}
Individual PDFs are named \code{`<project>_<tex-file>.pdf`}; the combined PDF is named \code{`<project>_combined.pdf`}; \code{`<project>`} is the folder name after \code{`docs-`}.

\section*{Compiling all docs}
\begin{DocCode}
  bash tools/render_all.sh [--fix] [--commit] [--do-not-retain-comments]
\end{DocCode}
This script regenerates renderable files in the normal \code{docs-*} folders (which contains a \code{manual/src/} folder)
Generated PDFs are written to each project's \code{pdf/} folder.
Combined PDFs are written to \code{current_pdfs/}.
By default, comments are retained when replacing a combined PDF in \code{current_pdfs/}.
The retention flow merges \code{comments/<pdf-stem>_comments.jsonl}, \code{comments/.base/<pdf-stem>_comments.jsonl}, and comments extracted from the previous \code{current_pdfs/<pdf-stem>.pdf}.
After the new current PDF is saved, the corresponding \code{comments/*.jsonl} and \code{comments/.base/*.jsonl} are refreshed from the saved PDF.

In case \code{--fix} is passed any compilation error will be tried to automatically be fixed by codex using this prompt
\begin{DocCode}
  Fix this compilation error:
  <Error>
\end{DocCode}
In case \code{--commit} is passed, for each folder a commit will be made (using the script \code{push_submodule_changes.sh})
In case \code{--do-not-retain-comments} is passed, the comments already existing in \code{current_pdfs/} are not retained.


\section*{General script prompt}
Whenever a script takes a long time there should be a progress bar!


\section*{Git submodule shortcuts}
This document gives shortcuts for pulling and pushing a Git repository together with all recursive submodules.
This is useful because each notes project is a separate Git repository, and some projects may contain their own submodules.
Our note system, and thus these commands, assume that the main repository and all submodules use the \code{main} branch, are up to date, and have an upstream remote configured.

\section*{Writing tex}
Each tex file follows first an overview before going more into details and implmentation.
Each \code{.tex} file should first provide an overview before going into implementation details.

\subsection*{Pull latest commits in the main repository and all submodules}
Use this when we want the main repository and every recursive submodule to move to the latest commit on \code{main}, and then record those updated submodule commits in their parent repositories.
\code{pull_update_submodules.sh}
\begin{DocCode}
msg="Update submodules"

git pull --ff-only
git submodule update --init --recursive

git submodule foreach --recursive '
git checkout main
git pull --ff-only
git submodule update --init --recursive
'

git submodule foreach --recursive '
git add -A
git diff --cached --quiet || git commit -m "$msg"
git push
'

git add -A
git diff --cached --quiet || git commit -m "$msg"
git push
\end{DocCode}

\subsection*{Add, commit, and push changes in all submodules}
Use this when we have edited files inside the main repository or inside any recursive
submodule, and we want to commit and push every repository separately.
Each submodule is committed and pushed first, then the parent repository records the updated
submodule commits and is pushed last.
\code{push_submodule_changes.sh}
\begin{DocCode}
msg="Update docs"

git submodule foreach --recursive '
git checkout main
git add -A
git diff --cached --quiet || git commit -m "$msg"
git push
'

git add -A
git diff --cached --quiet || git commit -m "$msg"
git push
\end{DocCode}

This assumes that the main repository and every submodule are on \code{main} and have an upstream remote configured.
The submodules are handled first because the parent repository only records the commit currently checked out inside each submodule.

\subsection*{Appendix: Comment JSONL schema}
\label{app:comment-jsonl-schema}
Each line in \code{comments/<pdf-stem>_comments.jsonl} is one JSON object representing one PDF annotation.
The same schema is used for \code{comments/.base/<pdf-stem>_comments.jsonl}.
\begin{itemize}
  \item \code{comment_index}: 1-based order within the extracted file.
  \item \code{source_pdf}: path of the PDF the comment was extracted from.
  \item \code{page}: 1-based PDF page number.
  \item \code{subtype}: PDF annotation subtype, such as \code{Text}, \code{Highlight}, \code{Underline}, \code{Squiggly}, or \code{StrikeOut}.
  \item \code{object_id}: PDF object id and generation, if available.
  \item \code{id}: PDF annotation \code{/NM} stable identifier, if available.
  \item \code{rect}: annotation rectangle.
  \item \code{quad_points}: highlighted text geometry, when present.
  \item \code{author}: annotation author.
  \item \code{modified}: PDF modified timestamp.
  \item \code{created}: PDF creation timestamp.
  \item \code{contents}: comment text.
  \item \code{color}: annotation color.
  \item \code{flags}: PDF annotation flags.
  \item \code{opacity}: annotation opacity.
  \item \code{name}: text annotation icon/name, mainly for sticky notes.
\end{itemize}
\code{comment_index}, \code{source_pdf}, and \code{object_id} are operational metadata and are ignored by semantic comment merging.
For example:
\begin{DocCode}
{"author":"Tommaso Bendinelli","color":[0.98,0.85,0.0],
 "comment_index":1,"contents":"Fix wording",
 "created":"D:20260706120000+02'00'","flags":28,
 "id":"example-comment-id","modified":"D:20260706120500+02'00'",
 "name":"Comment","object_id":[123,0],"opacity":1,
 "page":1,"quad_points":null,"rect":[72,720,93,741],
 "source_pdf":"/Users/tbe/Documents/docs/current_pdfs/template.pdf",
 "subtype":"Text"}
\end{DocCode}

\subsection*{Appendix: Information collection method}
\label{app:related-work-information-method}
We used Codex to query citation counts from Semantic Scholar and first-online dates from arXiv, specifically using the prompt:
\begin{DocCode}
  Find the number of citations for each article using Semantic Scholar, and the first time each article appeared online on arXiv, and update the document.
\end{DocCode}
Citation counts were retrieved with Codex using Semantic Scholar. 
The first-online metadata records (year and month) the earliest public online date or source date.
When a citation count or month/year could not be found from those sources, the value was completed with ChatGPT and marked with an asterisk. 
Dataset origin records whether the reviewed artifact introduces a new dataset, reuses a prior dataset, or is only described as a benchmark in the available notes.
Availability records a dataset URL only when the current document already identifies one; otherwise it uses \texttt{NaN} or notes that a public URL is not specified here.
Time-series length records the per-series or per-record length when available, notes source-dependent variation when lengths differ across source datasets, and uses ``Not applicable'' for non-time-series benchmarks.
Task goal records the primary evaluation target or dataset use, such as classification, forecasting, question answering, captioning, data analysis, or code-agent decision making.
Text fields record where natural language appears, such as question, instruction, context, system description, caption, report, annotation, answer, generated text, task specification, repository instruction, or bug report.


\end{document}
